\documentclass{article}

\usepackage[margin=0.85in]{geometry}
\usepackage{xcolor}
\usepackage{stepdiff}

\title{stepdiff Typed Color Examples}
\author{}
\date{}

\setlength{\parindent}{0pt}
\setlength{\parskip}{4pt}

\begin{document}

\maketitle
\vspace{-1.4em}

\noindent\small Single mode keeps one changed style; typed mode separates added and modified chunks; teaching mode also marks repeated operations such as \ensuremath{\lim} applied to both sides. The legend uses the actual visual styles.\normalsize\par

\section*{Single color mode}

\[
\begin{stepdiff}[display=notes, theme=soft, color-mode=single, frame=true]
\step{(x + 2)(x + 3)}
\step[reason={distribute}]{x(x + 3) + 2(x + 3)}
\step[reason={collect terms}, diff=all, tag=final]{x^2 + 5x + 6}
\end{stepdiff}
\]

\section*{Typed color mode}

\[
\begin{stepdiff}[display=notes, theme=soft, color-mode=typed, reason-style=muted, frame=true, legend=true]
\step{(x + 2)(x + 3)}
\step[reason={distribute}]{x(x + 3) + 2(x + 3)}
\step[reason={collect terms}, diff=all, tag=final]{x^2 + 5x + 6}
\end{stepdiff}
\]

\section*{Teaching mode: operation on both sides}

\[
\begin{stepdiff}[display=focus, theme=focus, color-mode=teaching, reason-style=badge, frame=true, legend=true]
\step{a_n \le b_n}
\step[reason={take limits}]{\lim a_n \le \lim b_n}
\step[reason={identify limits}, diff=all, tag=final]{L \le M}
\end{stepdiff}
\]

\section*{Teaching mode with final emphasis}

\[
\begin{stepdiff}[display=slide, theme=focus, color-mode=teaching, reason-style=badge]
\step{f(x) = x^2 + 3x}
\step[reason={differentiate}]{f'(x) = 2x + 3}
\step[reason={evaluate at x = 3}]{f'(3) = 2 \cdot 3 + 3}
\step[reason={final value}, diff=all, tag=final]{f'(3) = 9}
\end{stepdiff}
\]

\end{document}
